and a sheaf @@M0@@ of finite presentation on @@M1@@, such that @@M2@@ (@@M3@@ being the morphism @@M4@@). Since the @@M5@@ and the @@M6@@ are of finite presentation, we can, by EGA IV 11.2.6, find an index @@M7@@ such that @@M8@@ is flat over @@M9@@. All the hypotheses of 1.3 are therefore satisfied, which completes the proof.

\sgaStableTarget{sga6.e06.cor.1-4}\textbf{Corollary 1.4.} \uline{Under the hypotheses of} 1.3, \uline{there exists a graded} @@M10@@\uline{-Module} @@M11@@, \uline{of finite presentation}, \uline{and flat over} @@M12@@, \uline{such that} @@M13@@.

Let @@M14@@ be an integer such that the properties of 1.3 hold. If @@M15@@, we have seen that @@M16@@ is of finite presentation, and it is flat over @@M17@@, since the @@M18@@, for @@M19@@, are locally free on @@M20@@.

Moreover, @@M21@@ by EGA II 3.4.

Let us turn to part b) of the proof.

\sgaStableTarget{sga6.e06.lem.1-5-1}\textbf{Lemma 1.5.1.} \uline{Let} @@M22@@ \uline{be a ring}, @@M23@@ \uline{a positively graded} @@M24@@-\uline{algebra} \uline{such that} @@M25@@, @@M26@@ \uline{a graded} @@M27@@-\uline{module}. \uline{We suppose that} @@M28@@ \uline{is projective of finite type}, \uline{as an ungraded} @@M29@@-\uline{module}. \uline{Let} @@M30@@ \uline{be the sheaf of algebras on} @@M31@@ \uline{defined by} @@M32@@, @@M33@@ \uline{the} @@M34@@-\uline{Module associated with} @@M35@@; \uline{then} @@M36@@ \uline{is, locally on} @@M37@@, \uline{free of finite type over} @@M38@@, \uline{as a graded} @@M39@@-\uline{Module} (1.2).

Since @@M40@@ is of finite type, its degrees are bounded below. Moreover, @@M41@@ is projective of finite type over @@M42@@, and consequently has only finitely many nonzero homogeneous components, themselves projective of finite type over @@M43@@; by localizing on @@M44@@, we can therefore suppose that @@M45@@ is free of finite type over @@M46@@. Finally, since @@M47@@ is projective, @@M48@@. These properties imply that @@M49@@ is free of finite type over @@M50@@ as a graded module (see Bourbaki, Alg. ch. II §11 prop. 7), so @@M51@@ is free of finite type over @@M52@@, q.e.d.

